% RecSys Odyssey repository-backed lecture note.
% Add figures with: \coursefigure{figures/name.svg}{Caption}
\section{Recommendation is an intervention}

\begin{abstract}
A recommendation does not reveal a fixed preference; it changes the choice set.
\end{abstract}

\subsection{Concept}
Showing an item makes interaction possible and hides alternatives behind limited screen space. The observed response therefore combines user preference, presentation, position, timing and the policy that selected the item. A healthy analysis starts from impressions and eligible alternatives, not clicks alone. This distinction turns “what did users choose?” into the more honest question “what did users choose among what the system allowed them to see?”

\begin{equation*}
observed response = preference × exposure × presentation × context
\end{equation*}

\subsection{Key terms}
\begin{itemize}
\item \textbf{Choice set.} The alternatives available to a user at the moment of decision.
\item \textbf{Impression.} A logged exposure of an item in a known position and context.
\item \textbf{Intervention.} A system action that changes the user’s available experience.
\end{itemize}
